\section{Conclusion}

This article examined how much of the working day in the three largest occupations in community food, housing, and relief services is spent on tasks that AI can perform. Tasks that an LLM could accelerate without additional software account for less than an hour of the frontline working day. A larger share falls on tasks that would require software built for the purpose, and observed AI use reaches few tasks in these occupations. For organizations facing staff shortages, general-purpose AI tools are therefore likely to offer modest relief. There is greater potential in offloading worker burden via custom software built to support the task.

The article contributes a time-weighted account of AI exposure dedicated to the local community service relief sector. The article shows that measures based on counts of tasks versus time-weighted results may lead to different conclusions. Future research could test these estimates through direct time-use studies dedicated to the occupations at hand and also through usage data from more than one AI product. Studies of AI in the nonprofit sector would also benefit from recording which tasks AI is used for and how much working time those tasks occupy.
